\section{Evidence Coding, Adjudication, and Uncertainty}
\label{app:coding}

For each study, coding records the modeled object, temporal target, action semantics, transition mechanism, outcome, counterfactual support, planning role, and highest reported validation setting. The clinical facets include primary application, disease family, organ system, modality, and evaluated task. Public code, model, and data status are recorded separately from scientific capability because availability can change and does not establish validity.

A capability decision follows two stages. The existence gate first asks whether a clinically meaningful state and directed transition are implemented. The level decision then identifies the strongest task contract supported by the method and evaluation. Level~1 denotes state representation without directed transition; Level~2 denotes future prediction without a settable action; Level~3 requires a settable action that enters the transition and changes an evaluated future; Level~4 requires an explicit alternative-outcome identification or structural contract; and Level~5 requires operational use of a world model in action ranking, optimization, evaluation, or control. Ambiguity is retained in confidence and boundary notes rather than resolved from titles or self-description alone.

The review distinguishes unavailable evidence from negative evidence. A component is coded as not publicly evidenced when the paper and linked official artifacts do not describe it clearly enough to inspect. This designation does not imply that the component is absent from private implementation. Quantitative summaries include only canonical works with sufficient evidence for a capability decision; contextual and unresolved records are reported separately.

The manuscript uses Level~1--5 consistently. Legacy database values are transformed only at the display layer, so the naming change does not alter source records or quantitative results.

\begin{table}[t]
\caption{Capability distribution among implemented canonical works with sufficient evidence for adjudication ($n{=}217$). Each work is counted once; contextual, duplicate-version, unresolved, and non-transition records are excluded. Levels are non-cumulative: Level~5 does not certify Level~4.}
\label{tab:levels}
\centering
\small
\setlength{\tabcolsep}{4pt}
\renewcommand{\arraystretch}{1.08}
\begin{tabularx}{\linewidth}{@{}>{\raggedright\arraybackslash}p{2.55cm}Xrr@{}}
\toprule
\textbf{Capability category} & \textbf{Operational decision rule} & \textbf{$n$} & \textbf{\%} \\
\midrule
Level 1: current-state modeling & Current-state representation, calibration, or same-time synthesis; no directed future target. & 38 & 17.5 \\
Level 2: future prediction & Directed future state or observation without a settable action entering the transition. & 74 & 34.1 \\
Level 3: action-conditioned simulation & A settable action changes an evaluated, meaningful post-action future. & 35 & 16.1 \\
Level 4: counterfactual outcome reasoning & Alternative-intervention outcomes are estimated with an explicit causal identification contract. & 17 & 7.8 \\
Level 5: planning/control & A world model is operationally used to rank, select, optimize, or control actions or policies. & 53 & 24.4 \\
\bottomrule
\end{tabularx}
\end{table}

\begin{table}[t]
\caption{Primary clinical-application distribution among implemented clinical works with sufficient evidence for adjudication ($n{=}192$). Each canonical work is counted once. General-domain anchors and non-clinical biomedical context papers are excluded from this denominator.}
\label{tab:tracks}
\centering
\small
\setlength{\tabcolsep}{6pt}
\renewcommand{\arraystretch}{1.08}
\begin{tabularx}{\linewidth}{@{}Xrr@{}}
\toprule
\textbf{Primary clinical application} & \textbf{$n$} & \textbf{\%} \\
\midrule
Counterfactuals and treatment planning & 41 & 21.4 \\
EHR and patient trajectories & 39 & 20.3 \\
Longitudinal disease progression & 30 & 15.6 \\
Medical imaging representation & 27 & 14.1 \\
Surgical, robotic, and physiology models & 27 & 14.1 \\
Patient and health-system digital twins & 23 & 12.0 \\
Other medical settings & 4 & 2.1 \\
Population-health policy simulation & 1 & 0.5 \\
\midrule
\textbf{Total} & \textbf{192} & \textbf{100.0} \\
\bottomrule
\end{tabularx}
\end{table}

\begin{table}[t]
\caption{Validation ladder for medical world-model claims. The rungs are evidence types rather than a score: the required rung depends on the intended use, and evidence at one rung does not substitute for causal identification or action support at another.}
\label{tab:validation-pyramid}
\centering
\small
\setlength{\tabcolsep}{4pt}
\renewcommand{\arraystretch}{1.08}
\begin{tabularx}{\linewidth}{@{}p{2.35cm}p{3.35cm}Xp{3.25cm}@{}}
\toprule
\textbf{Evidence rung} & \textbf{What it can establish} & \textbf{Minimum useful evidence} & \textbf{Typical overclaim} \\
\midrule
Internal retrospective & Performance on held-out observations from the same data regime. & Patient-level splits, leakage controls, horizon-specific error, and relevant static baselines. & Generation quality is treated as decision validity. \\
External or multicenter & Robustness to site, scanner, coding, workflow, or population shift. & Prespecified external cohorts and subgroup-specific failure reporting. & One external mean score is treated as universal transportability. \\
Expert or clinician review & Clinical coherence and recognizable failure modes. & Blinded, task-specific ratings with agreement and explicit rejection criteria. & Plausibility substitutes for outcome or causal evidence. \\
Counterfactual or policy & Behavior under alternative actions within stated assumptions and support. & Explicit estimand, overlap diagnostics, sensitivity analysis, and trial-, simulator-, or semi-synthetic anchors. & Benchmark counterfactual accuracy is treated as patient-level truth. \\
Prospective workflow or trial & Effect on a real process, decision, or patient-relevant endpoint. & Registered protocol, monitoring, predefined endpoints, and appropriate comparator. & Retrospective gains are assumed to transfer to deployment. \\
\bottomrule
\end{tabularx}
\end{table}

